\documentclass[10pt,a4paper]{amsart}
\usepackage[margin=19mm]{geometry}
\usepackage{amsmath,amssymb,mathtools}
\usepackage[scr=rsfs,cal=euler]{mathalfa}
\usepackage{xfrac}
\usepackage[unicode,hidelinks,bookmarks=false]{hyperref}
\newcommand{\ie}{i.e.\ }
\newcommand{\cf}{cf.\ }
\newcommand{\resp}{respectively\ }
\newcommand{\Subsection}{\subsection}
\providecommand{\nobreakdash}{\nobreak\hbox{-}}
\setlength{\emergencystretch}{2em}
\multlinegap0pt
\usepackage{bm}
\usepackage{bbm}
\usepackage{dsfont}
\usepackage{xspace}
\usepackage{cancel}
\usepackage{mathtools}
\usepackage{cleveref}
\usepackage{amsthm}
\makeatletter
\@ifundefined{theorem}{\newtheorem{theorem}{Theorem}}{}
\@ifundefined{lemma}{\newtheorem{lemma}[theorem]{Lemma}}{}
\makeatother
\title[Latin squares with three disjoint subsquares of the same ord]{Latin squares with three disjoint subsquares of the same order}
\author{Tara Kemp, James G. Lefevre}
\date{Source: arXiv:2510.00364v1 (CC BY 4.0)}
\begin{document}
\maketitle

\noindent\emph{Source and licence.} Latin squares with three disjoint subsquares of the same order. Version arXiv:2510.00364v1 (CC BY 4.0), distributed under the Creative Commons Attribution 4.0 International licence, as recorded in the arXiv licence metadata for this exact version and shown on the abstract page https://arxiv.org/abs/2510.00364v1. Full rights notice: licenses/arxiv-2510.00364-CC-BY-4.0.txt.\par\smallskip

\noindent\textit{Editorial scope.} Counted unit: the Lemma whose source label is \texttt{even\_r\_from\_circulant} in the article (source0.tex lines 297--305, source label \texttt{even\_r\_from\_circulant}), delivered together with the article's own complete original proof of it (source0.tex lines 306--365, one \texttt{proof} environment of 60 lines). No mathematics is written, repaired or re-derived: statement and proof are the article's own text line for line. Required context retained uncounted: circulant\_construction (lines 222--232). Every definition, display and label of those lines is retained unchanged. Omitted from this package and not redistributed: 1--221, 233--296, 366--547. Nothing inside a retained excerpt is omitted or altered: every retained body is delivered exactly as the article's lines are written, so no omission receipt of any kind accompanies this package. Citations: the retained excerpts carry no citation command at all, so no bibliography entry of the article is transcribed and no reference is redistributed. Reference adaptation: none was needed and none was made; every reference inside the delivered unit resolves inside it, so no unresolved reference or citation is produced. Printed slips kept literal: no printed word, symbol, hypothesis or number of the counted statement or of its proof was altered. Layout and class: the article's own class is \texttt{sn-jnl} with packages including graphicx, geometry, xcolor, multirow, multicol, amsmath, amsthm, mathtools, amssymb, caption, algorithm, algpseudocode, bm, titlesec; the wrapper is the lane's pinned \texttt{amsart} editorial wrapper at 10pt on A4, which restates the theorem-like environments the retained text needs and adds the article's own macro definitions that the retained text needs (none), with extra wrapper packages (bm, bbm, dsfont, xspace, cancel, mathtools, cleveref). The delivered numbering is the unit's own. Assets: no retained span contains a figure environment, a graphics-inclusion command, a PDF-page inclusion, a table or a TikZ picture; no image file is copied into this package, the package declares no asset dependency and no image file is redistributed with it. Licence: version arXiv:2510.00364v1 of this article is distributed under the Creative Commons Attribution 4.0 International licence, as recorded in the arXiv licence metadata for this exact version and shown on the abstract page https://arxiv.org/abs/2510.00364v1. A full rights notice is delivered at licenses/arxiv-2510.00364-CC-BY-4.0.txt. Characters: every retained excerpt is pure ASCII, so the delivered engine, XeLaTeX with its default Latin Modern text font, reports no missing character.
\begin{lemma}
\label{circulant_construction}
    Let $P=(h_1h_2h_3\dots h_k)$, where $n=\sum_{i=1}^k h_i$, $t=n-h_1$ is odd, and $h_i\ge h_{i+1}$ for $2\le i \le k-1$. Further suppose that $h_2 \le (t+1)/4$ and $2h_2 \le h_1 \le t+1-2h_2$.
    Define a multiset $V = \sum_{i=2}^k h_i\{h_1+i-1\}$, so $|V|=t$.
    Then there exists an outline rectangle $O$ associated to $((1^n), (1^n), (1^{h_1}h_2h_3\dots h_k))$ with the following properties:
    \begin{enumerate}
        \item For any $x,y \in P[1]$, then $O(x,y)\in P[1]$. For any
        $i \in [k]\setminus \{1\}$ and $x,y \in P[i]$, then $O(x,y)=\{h_1+i-1\}$. Thus $O$ lifts to a realization of $P$ in normal form.
        \item We have triple sets $T_i=\{(x_{i,j},y_{i,j},z_{i,j}) \mid j\in [t]\}$, for all $i\in [h_1-2h_2]$, such that: $\{x_{i,j} \mid j\in [t]\}=\{y_{i,j} \mid j\in [t]\} =  [n]\setminus [h_1]$; $\{z_{i,j} \mid j\in [t]\}=V$; the pairs $(x_{i,j},y_{i,j})$ are distinct for all $i\in [h_1-2h_2]$ and $j \in [t]$; and for any $i\in [h_1-2h_2]$ and $j \in [t]$, $O(x_{i,j},y_{i,j})=\{i\}$ and $O(x_{i,j},i)=O(i,y_{i,j})=\{z_{i,j}\}$.
    \end{enumerate}
\end{lemma}

\begin{lemma}
\label{even_r_from_circulant}
    Let $P=(h_1h_2h_3h_4\dots h_k)$, where $h_1=h_2=h_3\ge 2$, $h_4 \ge \dots \ge h_k$, $h_1(h_1-1)\geq 2(h_k-1)$, $r=\sum_{i=4}^k h_i$ is even, $h_4 \le (r-2)/4$, and 
    $2h_4+2 \le 3h_1+1 \le r+1-2h_4$.
    
    Then there exists an outline rectangle $O$ associated to $(P,P,P)$ such that $O(i,i)$ contains $h_i^2$ copies of $i$ for each $i\in [k]$, and additionally $O(1,2)$, $O(2,3)$ and $O(3,1)$ contain $h_1^2$ copies of 3, 1, and 2 respectively, while 
    $O(1,3)$, $O(2,1)$ and $O(3,2)$ contain at least $h_1(h_1-1)-2(h_k-1)$ copies of 2, 3, and 1 respectively.

\end{lemma}

\begin{proof}
    Let $n=r+3h_1$.
By \Cref{circulant_construction} (with $t=r-1$), we have an outline rectangle $O_1$ associated to $((1^n), (1^n), (1^{3h_1+1}h_4h_5\dots (h_k-1)))$ which lifts to a realization of $((3h_1+1)h_4\dots (h_k-1))$ in normal form. The rows and columns $[3h_1+1]$ form a subsquare on symbols $[3h_1+1]$, and for $i\in [k-1] \setminus [3]$ the rows and columns $P[i]+1$ are filled with the symbol $3h_1+i-2$, while the rows and columns $(P[k]+1)\setminus \{n+1\}$ are filled with the symbol $3h_1+k-2$.
Further, since $(3h_1+1)-2h_4 \ge 2$, then there exist triple sets $T_1, T_2$ of size $r-1$ such that for $i\in [2]$ the following hold: (1) for any $x\in [n]\setminus [3h_1+1]$, there is exactly one triple $(x,y,z) \in T_i$ and one triple $(x,y,z) \in T_i$; (2) $\{z \mid (x,y,z) \in T_i$ is a multiset equal to $V= \left(\sum_{j=4}^{k} h_j \{3h_1+j-2\}\right)-\{3h_1+k-2\}$; and for any $(x,y,z) \in T_i$, cell $(x,y)$ contains $i$ and cells $(y,i)$ and $(i,x)$ contain $z$. 

    Let $P' = (h_1^31^r)$. We form an outline rectangle $O_2$ associated to $(P',P',(h_1^3 h_4h_5\dots h_{k-1}(h_k-1)1)$ by applying the following amalgamation and relabelling to $O_1$: rows, columns and symbols $[3h_1]$ are amalgamated into 3 groups of size $h_1$, labelled 1,2,3, such that the original values 1 and 2 are joined to groups 1 and 2 respectively. The row and column values $[n] \setminus [3h_1+1]$ are each reduced by $3h_1-2$, mapping to the range $4$ to $r+2$, while $3h_1+1$ becomes $r+3$. Similarly, the symbol values $[k+3h_1-2] \setminus [3h_1+1]$ are each reduced by $3h_1-2$, mapping to the range $4$ to $k$, while $3h_1+1$ becomes $k+1$.
    Finally, rows 2 and 3 are swapped, columns 1 and 3 are swapped, and symbols 1 and 2 are swapped. 
    The cells $O_2(i,j)$, $i,j \in \{1,2,3,r+3\}$, contain only elements of $\{1,2,3,k+1\}$, and thus represent a subsquare. So we can make the substitution 
\newline

    $\begin{array}{c|c|c|c|c|}
O_2(i,j) & j=1 & j=2 & j=3 & j=r+3 \\ 
\hline
i=1 & h_1^2 \{1\} & h_1^2\{3\} & h_1(h_1-1)\{2\} & h_1\{2\} \\
&&&+h_1\{k+1\}& \\
\hline
i=2 & h_1(h_1-1)\{3\} & h_1^2\{2\} & h_1^2 \{1\} & h_1\{3\} \\
&+h_1\{k+1\}&&& \\
\hline
i=3 & h_1^2\{2\} & h_1(h_1-1)\{1\} & h_1^2\{3\} & h_1\{1\} \\
&&+h_1\{k+1\}&& \\
\hline
i=r+3 & h_1\{3\} & h_1\{1\} & h_1\{2\} & \{k+1\} \\
\hline
\end{array}$
\newline

    We now form an outline rectangle $O_3$ associated to $(P',P',P)$ by further amalgamating the symbols $k$ and $k+1$ in $O_2$, with the new symbol labelled $k$.
    Observe that for each $i \in [k]\setminus \{3\}$, and every $x,y\in (P[h_i]-3h_1+3)$, $O_3(x,y)=\{i\}$, with the exception of $i=k$ and exactly one of $x,y$ equal to $r+3$.
    To complete the construction of $O$, we therefore identify trades that replace the multisets $O_3(x,r+3)$ and $O_3(r+3,x)$ with $\{k\}$, for all $x \in (P[h_k]-3h_1+3) \setminus \{r+3\}$. 

    For $i=1,2$, let $T'_i = \{(y-3h_1+2,x-3h_1+2,z-3h_1+2) | (x,y,z) \in T_i\}$.
    Then for each $(x,y,z) \in T'_1$, we have $O_3(x,y)=\{2\}$ and $z \in O_3(1,y), O_3(x,3)$, and for each $(x,y,z) \in T'_2$, we have $O_3(x,y)=\{1\}$ and $z \in O_3(x,2), O_3(3,y)$.

    We identify four sequences of size $h_k-1$: for $i \in [h_k-1]$ let $a_{i} = r+3-i$, and let $b_i < r+3$, $c_i<r+4-h_k$ and $d_i$ be the values satisfying $O_3(r+3,b_i)=\{k\}$, $O_3(c_{i},a_{i})=\{k\}$, and $O_3(r+3,a_{i})=\{d_{i}\}$. Note that $b_i,c_i,d_i \in [r+2]\setminus [3]$. For $b_i$ and $d_i$, this follows from the fact that for each $x\in[3]$, $O_3(r+3,x)=O_2(r+3,x)=h_1{y}$ for some $y\in[3]$. For $c_i$, we observe that $O_3(x,a_i)=O_2(x,a_i)=\{k\}$ for all $x \in [r+2]\setminus [r+3-h_k]$. The cell $(c_i,a_i)$ contains the remaining copy of $k$ in column $a_i$ of $O_3$, which corresponds to the copy of $k+1$ in column $a_i$ of $O_2$. In the array above we see that rows $1,2,3,r+3$ of $O_2$ contain all available copies of $k+1$ in columns $1,2,3,r+3$, so $c_i\in [r+2]\setminus [3]$.
    Let $U = \{(y_j,x_j,z_j) \mid j \in [u]\}$ be the minimal subset of $T'_1$ such that $a_i, b_i \in \{x_j \mid j \in [u]\}$, $c_i \in \{y_j \mid j \in [u]\}$, $d_i \in \{z_j \mid j \in [u]\}$. Since $\{a_i \mid i \in [h_k-1]\}$ and $\{b_i \mid i \in [h_k-1]\}$ must be disjoint, $2(h_k-1) \le u=|U| \le 4(h_k-1)$. The trade is defined in the following table:
    \newline

    $\begin{array}{c|c|c|c}
    \text{cell} & \text{parameter} & \text{removed} & \text{added} \\
    \hline
    (y_j,x_j) & j\in [u] & \{2\} & \{z_j\} \\
    (y_j,3) & j\in [u] & \{z_j\} & \{2\} \\
    (1,x_j) & j\in [u] & \{z_j\} & \{2\} \\
    (r+3,a_i) & i\in [h_k-1] & \{d_i\} & \{k\} \\
    (r+3,b_i) & i\in [h_k-1] & \{k\} & \{2\} \\
    (r+3,3) && (h_k-1)\{2\} & \{d_i \mid i \in [h_k-1]\} \\
    (c_i,3) & i\in [h_k-1] & \{2\} & \{k\} \\
    (c_i,a_i) & i\in [h_k-1] & \{k\} & \{2\} \\
    (1,a_i) & i\in [h_k-1] & \{2\} & \{d_i\} \\
    (1,b_i) & i\in [h_k-1] & \{2\} & \{k\} \\
    (1,3) & & u\{2\} & \{z_j \mid j\in [u] \} \\
    (1,3) & & \{d_i \mid i\in [h_k-1] \} & (h_k-1)\{2\}\\
    (1,3) & & (h_k-1)\{k\} & (h_k-1)\{2\} \\
\end{array}$
\newline

This trade replaces $O_3(r+3,i)$ with $\{k\}$ for $i\in (P[h_k]-3h_1+3) \setminus \{r+3\}$. A similar trade replaces $O_3(i,r+3)$ with $\{k\}$: rows and columns are reversed, and $T_2$, row 3, column 2 and symbol 1 are used in place of $T_1$, row 1, column 3, and symbol 2. These trades are disjoint from each other and from the remainder of the required subsquares. We apply both of these trades and then amalgamate rows and columns to give $O$ associated to $(P,P,P)$; each row and column set $P[i]-3h_1+3$ is mapped to $i$, for $i\ge 4$. $O_3(1,3)$ contains $h_1(h_1-1)$ copies of $2$, which is reduced by $u-2(h_k-1)$ in the trade. Thus $O(1,3)$ contains $h_1(h_1-1)+2(h_k-1)-u \ge h_1(h_1-1)-2(h_k-1)$ copies of $2$, and similarly $O(3,2)$ contains at least $h_1(h_1-1)-2(h_k-1)$ copies of $1$.

\end{proof}

\end{document}
